% Example "burgers": the problem, the shock and rarefaction in space-time, the
% overshoots of the methods, and the total-variation test. Every number comes
% from data/, written by examples/burgers/main.f90: run `fpm run --example
% burgers` from the project root, then build this file with lualatex in
% examples/burgers/.
\documentclass[11pt,a4paper]{article}
\usepackage[margin=2.2cm]{geometry}
\usepackage{amsmath,booktabs,array}
\usepackage{pgfplots}
\usepgfplotslibrary{groupplots,colormaps,fillbetween}
\pgfplotsset{compat=1.18}
\pagestyle{empty}

\colorlet{m-euler}{black}
\colorlet{m-order2}{orange!85!black}
\colorlet{m-order3}{green!50!black}
\colorlet{m-rk4}{blue!80!black}
\colorlet{m-dp54}{red!75!black}
% Diverging colour map for u: blue (u = -0.5), white (u = 0.25), red (u = 1).
\pgfplotsset{colormap={updown}{rgb255(0cm)=(33,102,172) rgb255(1cm)=(146,197,222)
  rgb255(2cm)=(247,247,247) rgb255(3cm)=(244,165,130) rgb255(4cm)=(178,24,43)}}

\begin{document}
\section*{Example: Burgers' equation and strong stability}

\subsection*{1. Problem}
The inviscid Burgers equation, the simplest model of a nonlinear wave that
steepens into a shock, with a square wave as initial value:
\begin{equation*}
  u_t + \bigl(\tfrac12 u^2\bigr)_x = 0, \quad 0 \le x < 1 \text{ periodic}, \qquad
  u(x, 0) = \begin{cases} 1, & 0.25 < x < 0.6,\\ -0.5, & \text{elsewhere.}\end{cases}
\end{equation*}
The jump at $x = 0.6$ is a shock moving at the Rankine--Hugoniot speed
$\tfrac12(1 - 0.5) = \tfrac14$; the jump at $x = 0.25$ opens into a
rarefaction fan between the characteristics of speeds $-0.5$ and $1$. Space
is discretized by a TVD scheme (MUSCL reconstruction, minmod limiter, Rusanov
flux): with forward Euler it never increases the total variation
$\mathrm{TV}(u) = \sum_i |u_{i+1} - u_i|$, and so never creates new extrema,
for CFL $= \max|u|\,\Delta t/\Delta x \le \tfrac12$. The question is which
Runge--Kutta methods keep that property.
\begin{center}
\small
% Written by examples/burgers/main.f90: do not edit.
\begin{tabular}{l>{\raggedright\arraybackslash}p{12.4cm}}
\toprule
domain & $0 \le x < 1$, periodic; $200$ cells, $\Delta x = 1/200$ \\
initial condition & $u(x, 0) = 1$ on $0.25 < x < 0.6$, $-0.5$ elsewhere (a square wave) \\
interval & $0 \le t \le 0.4$ \\
space & MUSCL, minmod limiter, Rusanov flux: TVD with forward Euler for CFL $\le 1/2$ \\
time step & $\Delta t = \mathrm{CFL}\,\Delta x / \max|u_0|$, rounded down to land on $t = 0.4$ \\
reference & \texttt{rk4}, CFL $= 0.05$ \\
CFL sweep & $0.05, 0.10, \dots, 1.2$; largest $\mathrm{TV}(u_{n+1}) - \mathrm{TV}(u_n)$ and largest new extremum $\max(\max u - 1, -0.5 - \min u)$ over the run; TVD up to the first CFL with an increase above $10^{-12}$ \\
figures & solutions at CFL $= 0.6$ \\
methods & \texttt{rk3ws}, \texttt{rk4}, \texttt{dp54} \\
interface & plain arrays, \texttt{rk\_array\_scheme\_t}, 200 unknowns \\
\bottomrule
\end{tabular}

\end{center}

\newpage
\subsection*{2. Solutions}
\begin{center}
\begin{tikzpicture}
  % ---- (a) space-time ----
  \begin{axis}[
      name=spacetime,                                      % anchor for the panels on the right
      width=8.0cm, height=11cm,                            % tall map
      view={0}{90}, xmin=0, xmax=1, ymin=0, ymax=0.4,      % seen from above: x across, t up
      xlabel={$x$}, ylabel={$t$}, ytick={0,0.1,0.2,0.3,0.4}, % space-time
      title={(a) $u(x, t)$, reference},                    % panel title
      colormap name=updown, point meta min=-0.5, point meta max=1, % fixed range of u
      colorbar horizontal, colorbar style={at={(0.5,-0.13)}, anchor=north, xlabel={$u$}, height=0.2cm, width=6cm},
      enlargelimits=false, axis on top,                    % the map fills the axes
    ]
    \addplot3[surf, shader=interp, mesh/ordering=x varies] table[x=x, y=t, z=u] {data/spacetime.dat}; % the solution
    \addplot[black, thick, dashed, domain=0:0.4] ({0.6 + 0.25*x}, {x});        % shock path x = 0.6 + t/4
    \addplot[black, thin, dotted, domain=0:0.4] ({0.25 - 0.5*x}, {x});         % fan edge, speed -0.5
    \addplot[black, thin, dotted, domain=0:0.4] ({0.25 + x}, {x});             % fan edge, speed 1
    \node[font=\scriptsize, fill=white, fill opacity=0.8, text opacity=1, anchor=south] at (axis cs:0.69,0.30) {shock};
    \node[font=\scriptsize, fill=white, fill opacity=0.8, text opacity=1, anchor=south] at (axis cs:0.30,0.30) {fan};
  \end{axis}
  % ---- (b) final profiles ----
  \begin{axis}[
      name=final,                                          % anchor for (c)
      at={($(spacetime.north east)+(1.5cm,0cm)$)}, anchor=north west, % right of (a)
      width=6.6cm, height=5.6cm,                           % panel
      xmin=0, xmax=1, ymin=-0.65, ymax=1.15,               % the profile
      xlabel={$x$}, ylabel={$u$},                          % solution against space
      title={(b) $t = 0.4$, CFL $= 0.6$},                  % panel title
      axis background/.style={fill=black!2},               % very light background
      legend style={font=\tiny, at={(0.02,0.98)}, anchor=north west, fill=white, fill opacity=0.85, text opacity=1, draw=black!30},
    ]
    \addplot[black!25, line width=3pt] table[x=x, y=u] {data/reference.dat};           % reference
    \addlegendentry{reference}
    \addplot[m-dp54, thick, dashed] table[x=x, y=u] {data/final_dp54.dat};
    \addlegendentry{\texttt{dp54}}
  \end{axis}
  % ---- (c) the first steps, zoomed at the shock ----
  \begin{axis}[
      at={($(final.south west)+(0,-1.6cm)$)}, anchor=north west, % below (b)
      width=6.6cm, height=6cm,                             % panel
      xmin=0.575, xmax=0.625, ymin=-0.62, ymax=1.1,        % zoom on the jump at x = 0.6
      xlabel={$x$}, ylabel={$u$},                          % solution against space
      title={(c) after two steps, CFL $= 0.6$},            % panel title
      axis background/.style={fill=black!2},               % very light background
      legend style={font=\tiny, at={(0.98,0.98)}, anchor=north east, fill=white, fill opacity=0.85, text opacity=1, draw=black!30},
    ]
    \fill[red!8] (axis cs:0.575,1) rectangle (axis cs:0.625,1.1);                      % above the initial maximum
    \fill[red!8] (axis cs:0.575,-0.62) rectangle (axis cs:0.625,-0.5);                 % below the initial minimum
    \addplot[black, thin, densely dotted, domain=0.575:0.625, forget plot] {1};         % initial maximum
    \addplot[black, thin, densely dotted, domain=0.575:0.625, forget plot] {-0.5};      % initial minimum
    \addplot[m-dp54, dashed, mark=pentagon, mark size=1.6pt] table[x=x, y=u] {data/early_dp54.dat};
    \addlegendentry{\texttt{dp54}}
  \end{axis}
\end{tikzpicture}
\end{center}
(a) The reference solution: the shock (dashed, the Rankine--Hugoniot path
$x = 0.6 + t/4$) and the rarefaction fan (between the dotted characteristics
$x = 0.25 - t/2$ and $x = 0.25 + t$) appear exactly where the theory puts
them. (b) At $t = 0.4$ the methods agree with the reference at this scale.
(c) The difference is at the start, while the jump is still sharp: two steps
in, \texttt{rk3ws} and \texttt{rk4} stay between the initial bounds $-0.5$ and
$1$, while \texttt{dp54} overshoots to $1.006$ and $-0.506$ (shaded: new
extrema).

\newpage
\subsection*{3. Test: total variation and new extrema}

\paragraph{Prediction.} A Runge--Kutta method that is a convex combination of
forward Euler steps, with SSP coefficient $C$, keeps any property forward
Euler has for steps up to $C$ times larger (Shu and Osher; Gottlieb, Ketcheson
and Shu). An SSP method such as SSP-RK3 ($C = 1$), which a program can give
the library as a tableau, would be guaranteed TVD up to CFL $\tfrac12$ (the
shaded range). The three methods of the library, \texttt{rk3ws}, \texttt{rk4}
and \texttt{dp54}, have $C = 0$: no guarantee at any step; they may still be
TVD in practice. The guarantee is sufficient, not necessary.

\paragraph{Result.}
\begin{center}
\footnotesize
\setlength{\tabcolsep}{4pt}
% Written by examples/burgers/main.f90: do not edit.
\begin{tabular}{lcccccc}
\toprule
& & & \multicolumn{2}{c}{largest TV increase per step} & \multicolumn{2}{c}{largest new extremum} \\
\cmidrule(lr){4-5}\cmidrule(lr){6-7}
name & SSP coeff. & TVD up to CFL & at CFL 0.45 & at CFL 0.60 & at CFL 0.45 & at CFL 0.60 \\
\midrule
\texttt{rk3ws} & 0 & 0.45 & $\le 10^{-12}$ & $3.22\times10^{-4}$ & $\le 10^{-12}$ & $1.61\times10^{-4}$ \\
\texttt{rk4} & 0 & 0.40 & $4.33\times10^{-11}$ & $\le 10^{-12}$ & $2.16\times10^{-11}$ & $\le 10^{-12}$ \\
\texttt{dp54} & 0 & 0.25 & $6.34\times10^{-4}$ & $5.62\times10^{-2}$ & $2.17\times10^{-4}$ & $1.63\times10^{-2}$ \\
\bottomrule
\end{tabular}

\end{center}

\begin{center}
\begin{tikzpicture}
  \begin{groupplot}[
      group style={group size=2 by 1, horizontal sep=1.8cm}, % two panels
      width=7.4cm, height=6.6cm,                           % size of each panel
      xmin=0, xmax=1.25, ymode=log, ymin=1e-16, ymax=10,   % CFL sweep, log of the increase
      xlabel={CFL}, grid=major, grid style={black!8},      % shared axes
      axis background/.style={fill=black!2},               % very light background
      every axis plot/.append style={thick, mark size=1.3pt}, % same weight for all curves
    ]
    % ---- (i) TV increase ----
    \nextgroupplot[
      title={(i) largest TV increase per step},            % panel title
      legend to name=bglegend, legend columns=7,           % legend below the panels
      legend style={font=\small, /tikz/every even column/.append style={column sep=0.25cm}},
    ]
    \fill[green!8] (axis cs:0,1e-16) rectangle (axis cs:0.5,10);                        % the SSP guarantee
    \node[font=\scriptsize, text=green!40!black, anchor=north, align=center] at (axis cs:0.25,1) {guaranteed TVD\\for $C = 1$};
    \addplot[black!50, thin, dashed, forget plot] coordinates {(0,1e-12) (1.25,1e-12)}; % the TVD tolerance
    \addplot[m-order3, dashed, mark=triangle] table[x=cfl, y=increase] {data/tv_rk3ws.dat};
    \addlegendentry{\texttt{rk3ws}}
    \addplot[m-rk4, dashed, mark=diamond]     table[x=cfl, y=increase] {data/tv_rk4.dat};
    \addlegendentry{\texttt{rk4}}
    \addplot[m-dp54, dashed, mark=pentagon]   table[x=cfl, y=increase] {data/tv_dp54.dat};
    \addlegendentry{\texttt{dp54}}
    % ---- (ii) new extrema ----
    \nextgroupplot[title={(ii) largest new extremum}]
    \fill[green!8] (axis cs:0,1e-16) rectangle (axis cs:0.5,10);                        % the SSP guarantee
    \addplot[m-order3, dashed, mark=triangle] table[x=cfl, y=extremum] {data/extrema_rk3ws.dat};
    \addplot[m-rk4, dashed, mark=diamond]     table[x=cfl, y=extremum] {data/extrema_rk4.dat};
    \addplot[m-dp54, dashed, mark=pentagon]   table[x=cfl, y=extremum] {data/extrema_dp54.dat};
  \end{groupplot}
  \node[below=0.2cm] at (current bounding box.south)        % legend centred under the panels
    {\pgfplotslegendfromname{bglegend}};
\end{tikzpicture}
\end{center}
Without the guarantee, the three methods lose the property inside the range
an SSP method would keep, or at its edge: \texttt{dp54} from CFL $0.30$,
\texttt{rk4} from $0.45$ (by $4\times10^{-11}$), \texttt{rk3ws} from $0.50$;
increases at the level of round-off ($\approx 3\times10^{-15}$, below the
dashed tolerance) are not counted. \texttt{rk4} is not monotone in the CFL:
TVD again between $0.55$ and $0.70$.
\end{document}
